\documentclass[10pt,a4paper]{article}
\usepackage[T1]{fontenc}
\usepackage[utf8]{inputenc}
\usepackage{lmodern}
\usepackage{microtype}
\usepackage[margin=2.2cm]{geometry}
\usepackage{booktabs}
\usepackage{tabularx}
\usepackage{array}
\usepackage[hidelinks]{hyperref}
\newcommand{\code}[1]{\texttt{#1}}
\newcolumntype{L}{>{\raggedright\arraybackslash}X}

\title{P-System: consensusless payments across committee changes,\\
with rounds that end a split vote and a Lean core that runs in the node}
\author{Paolo Amodio}
\date{P-System \code{1.0.0}. Not submitted, not published, not peer reviewed.}

\begin{document}
\maketitle

\begin{abstract}
Payment systems built on Byzantine quorum certificates need no consensus when each account has one
owner, as FastPay and Astro show. Two problems remain: changing the committee, and a sender who signs
two versions of one spend, which can lock that spend for good. We describe a protocol that changes
committee with no consensus and ends a split vote with rounds per key, and its implementation. Its
model is in Lean~4: safety over a chain of committees, and progress within an epoch under partial
synchrony. The decisions of the running node (the vote lock, the execution of a certificate, the steps of
a round) are the definitions of the model compiled to a shared library, with theorems that what the
Python around them passes is enough. That Python is tested and not proved. We report fault campaigns on
four virtual machines of one host, a liveness failure they found, and its correction.
Every theorem named here is in \code{formal/lean/PSystem}, every number in a file of the repository,
and \code{tests/test\_documents.py} checks both.
\end{abstract}

\section{The protocol}

A committee has $n = 3f+1$ members, at most $f$ of them faulty. Each key $(\mathit{sender},
\mathit{nonce})$ is decided on its own; there is no global order.

\paragraph{Round 0.} An honest member votes for the first version of a key it receives, durably, and
never for another. The votes of all $n$ members decide that version at once. $2f+1$ of them are a
\emph{prepare} certificate: each member that holds one precommits it, and $2f+1$ precommits decide.
A prepare alone is never executed: with $3f+1$ members no rule on $2f+1$ votes alone can be recovered
safely (\code{split\_lock\_ambiguous\_general}).

\paragraph{Rounds.} A member with an open key moves to round $r+1$ at a timeout (2~s, doubling up to
60~s) and sends a signed report to all: its vote of round 0 and the highest prepare it holds. The leader
of round $r$ is member $(k + r) \bmod n$, $k$ from the hash of the key. With $2f+1$ reports it proposes
the value of the highest prepare among them, or else one that $f+1$ voted in round 0, or else any version
its sender signed. A member prevotes a proposal only if it comes from the leader and the reports justify
it; $2f+1$ prevotes are a prepare of the round, and the precommits follow. Until a key is decided, a
member sends again, every 2~s, the messages it sent in its current round.

\paragraph{Epochs.} A threshold identity of governance signs an \emph{order} naming the next
committee. A member that receives it stops voting, waits up to 8~s for its open keys, and gives a
\emph{report}: every version it voted for in this epoch, in any round, with its sender's signature. It
reports for one order per epoch, for good. An order with $2f+1$ valid reports is a \emph{closure}; whoever
gathers one installs it, with no vote on it. Every version in the reports becomes a \emph{lock}: later
committees vote only for a version that agrees with the locks of its key.

\paragraph{Snapshots.} A state with the ids of the events behind it, signed by $2f+1$ members, each
of which derived it. A node can restart from one, prune what it covers, and apply a floor under which it
no longer votes.

\section{What is proved}

All of this is about the model. There is no \code{sorry}, \code{admit}, \code{axiom}, \code{opaque}
or \code{native\_decide}, and \code{formal/lean/AuditPSystem.lean} prints only \code{propext},
\code{Quot.sound} and \code{Classical.choice} for each of 269 theorems in 27 modules. Signatures are
parameters, and \code{Unforgeable} is a hypothesis about a run: nothing here is about Ed25519 or SHA-256.

\begin{center}\small
\begin{tabularx}{\textwidth}{@{}p{0.36\textwidth}L@{}}
\toprule
Theorem & Statement \\
\midrule
\code{decisions\_agree}, \code{grounds\_agree} & Two honest members never decide different versions of a key, whatever the network did, through any number of rounds \\
\code{good\_round\_decides} & In a round after those the honest members are in, with an honest leader and the messages between honest members delivered, every honest member decides \\
\code{honest\_leader\_within\_n\_rounds} & With fewer than $n$ faulty members, any $n$ consecutive rounds have an honest leader \\
\code{voted\_in\_closure}, \code{no\_double\_spend\_next\_epoch} & A decided version is in every closure whose reporters list what they voted for in any round, and the next epoch decides the same \\
\code{no\_double\_spend}, \code{state\_is\_derive}, \code{convergence\_after\_exchange} & For one committee: honest nodes never execute conflicting events, a state is a function of the set of events, nodes that exchange logs agree \\
\code{no\_double\_spend\_across\_epochs}, \code{closure\_preserves\_certified} & Over a chain of committees, each with its own bound, two conflicting events never both get a certificate \\
\code{order\_unique}, \code{no\_double\_spend\_any\_closures} & Two orders cannot both get a quorum of reports, and members can install different closures of one epoch \\
\code{certified\_passes\_threshold}, \code{no\_double\_spend\_with\_unlock} & With the optional $n = 5f+1$, $q = 4f+1$, freeing a nonce with fewer than $2f+1$ reports creates no conflict \\
\code{accepted\_complete}, \code{old\_keys\_no\_double\_spend} & A certificate of an ended committee counts only for an event in the locks; with one version of each spend in the locks, this is safe with stolen keys \\
\code{resume\_is\_derive}, \code{no\_double\_spend\_views} & A restart from a snapshot reaches the derived state; pruning locks under a floor keeps safety \\
\code{step\_certificate}, \code{executes\_local}, \code{apply\_local}, \code{voteAllowed\_local}, \code{execValue\_json} & The vote lock and the execution the node calls are those of the model, on the arguments Python passes (section~\ref{sec:link}) \\
\bottomrule
\end{tabularx}
\end{center}

Across epochs, safety is about certification: the model of epochs has no execution, which is proved
for one committee only. Each hypothesis of the epoch results has a counterexample that shows it is
needed, for instance \code{prune\_without\_floor\_lets\_conflict\_through} and
\code{order\_unique\_needs\_the\_bound}.

\section{Progress}

No deterministic protocol decides with one crash and no bound on delays (FLP), so progress needs a
hypothesis. \code{good\_round\_decides} states it as a good round: one after the rounds the honest members
are in, with an honest leader, in which their messages arrive. What happened before is arbitrary. With
timeouts that double, honest members end up in such a round once delays are bounded, and
\code{honest\_leader\_within\_n\_rounds} gives the leader. The two theorems \code{split\_votes\_no\_quorum}
and \code{split\_then\_good\_round\_decides} are the case that used to stall: four members, votes split two
and two, no prepare in round 0, and a good round after it decides at every member.

The model does not say that the messages of a round arrive the first time they are sent, and the
campaigns showed why it matters (section~\ref{sec:vm}). The node now sends them again within the round.

What progress does not cover: a split still open at the end of an epoch, after the wait of 8~s, is
reported with both versions and stays locked in the next epoch. Waiting without a bound could stop the
change of epoch, so we chose the bound.

\section{Python and Lean}\label{sec:link}

A proof that Python code is equivalent to a model needs a formal semantics of Python, and there is no
verified translation from Python to Lean. We did not attempt it. Instead the decisions are the model's
code. \code{lake build PSystem:shared} compiles the Lean definitions to a library that the node loads
(\code{network/quorum/lean\_core.py}): \code{Quorum.voteAllowed} for the vote lock, \code{executes} and
\code{ApplyEvent} for the execution of a certificate, and \code{Rounds.step}, \code{choose},
\code{justified} and \code{leader} for the rounds.

What can be proved about the boundary is proved in \code{CoreExec.lean} and \code{Core.lean}:
\code{step\_certificate} says the execution in \code{Quorum.step} is \code{executes} with each signature
replaced by the result of checking it; \code{executes\_local}, \code{apply\_local} and
\code{apply\_elsewhere} that the rows of the sender and the receiver, and the executed ids among the event
and its parents, are all the state the answer depends on; \code{voteAllowed\_local} that the lock depends
only on the votes with the key or the id of the event, which is what the store selects; and
\code{execValue\_json}, \code{voteValue\_json} that on the JSON encoding of any argument the exported
function returns the encoding of the model's answer.

What remains trusted is listed in \code{docs/PROOF-OBLIGATIONS.md}: the Python that builds the
arguments and writes the answer back, \code{Json.parse} and \code{Json.compress}, \code{ctypes}, the Lean
compiler and runtime, and what the model leaves out (Ed25519, SQLite, the network, the driver around the
decisions, the rule on keys of ended committees). It is tested: differential tests compare the node with
the Lean models after each input, on adversarial random inputs, for one replica, through the node, on real
TLS, across epochs and snapshots with restarts; mutation tests check that each rule of the remaining Python
is noticed when it is broken.

\section{Findings}

\paragraph{A closure needs no agreement.} Our first design voted on one closure per epoch. Two proposers
could split those votes so that no closure ever formed, after voting on events had stopped. Since
\code{closure\_preserves\_certified} holds for any quorum of reports, the votes are not needed; what must
be unique is the next committee, and a lock on the order gives that.

\paragraph{A prepare is not a decision.} Moving from round-0 certificates of $2f+1$ to rounds, the
differential tests found two gaps in the node: a prepare received by synchronization was taken as a
decision, and could halt a node on the real decision of the other version; and a prepare received was
not sent on, so a gatherer that showed it to one member left the others without a precommit until a
round ended.

\paragraph{Old keys.} With the keys of an ended committee, stolen after its closure, our implementation
executed a transfer its sender never signed. A node now takes a certificate of an ended committee only
for an event in its locks.

\paragraph{A message sent once.} A campaign on virtual machines found honest transfers not final two
minutes after a long partition healed (section~\ref{sec:vm}).

\section{Fault campaigns on virtual machines}\label{sec:vm}

VMRESULTS

\section{Related work}

What each work is reported to do, from its paper; none has been run against this system.

\begin{center}\small
\begin{tabularx}{\textwidth}{@{}p{0.2\textwidth}LLL@{}}
\toprule
& Consensus & Committee change & Equivocated spend \\
\midrule
FastPay & None & Fixed committee & Locked \\
Astro & None, reliable broadcast & Not handled & Not its subject \\
Kuznetsov, Tonkikh & None on the configuration & Its subject, forward-secure signatures & Not its subject \\
Sui Lutris & For shared objects and at the epoch boundary & Through consensus & Unlocked at the end of an epoch \\
Bythos & Verifies broadcast protocols in Coq, with code extraction & Not covered & Not covered \\
This work & None & By a threshold identity & Decided by rounds within an epoch; locked if still split at its end \\
\bottomrule
\end{tabularx}
\end{center}

Bythos proves more about its protocols and extracts running code from the proofs. Here the decisions are
compiled from the model, but the code around them is written by hand and tested.

\section{Limits}

\begin{itemize}
\item The Python around the Lean core is tested, not proved, and so is everything the model leaves out.
\item Progress needs partial synchrony; a split open at the end of an epoch stays locked.
\item Governance chooses the committee; more than $f$ faulty members can fork, and a governance that signs
two orders stops an epoch for good.
\item The campaigns ran four virtual machines on one host: never separate machines, independent operators,
an external audit or a pilot.
\item Only single-owner state. Every member sees every event.
\item One author, working with an AI model that wrote much of the code and the proofs. The proofs are
checked by Lean; the rest has been checked by no independent person.
\end{itemize}

\begin{thebibliography}{9}\small
\bibitem{guerraoui} R.~Guerraoui, P.~Kuznetsov, M.~Monti, M.~Pavlovic, D.-A.~Seredinschi. The Consensus Number of a Cryptocurrency. PODC 2019. arXiv:1906.05574.
\bibitem{fastpay} M.~Baudet, G.~Danezis, A.~Sonnino. FastPay: High-Performance Byzantine Fault Tolerant Settlement. AFT 2020. arXiv:2003.11506.
\bibitem{astro} D.~Collins et al. Online Payments by Merely Broadcasting Messages. DSN 2020. arXiv:2004.13184.
\bibitem{kt} P.~Kuznetsov, A.~Tonkikh. Asynchronous Reconfiguration with Byzantine Failures. DISC 2020. doi:10.4230/LIPIcs.DISC.2020.27.
\bibitem{sui} S.~Blackshear et al. Sui Lutris: A Blockchain Combining Broadcast and Consensus. CCS 2024. arXiv:2310.18042.
\bibitem{bythos} Q.~Zhao, G.~P\^irlea, K.~Grzeszkiewicz, S.~Gilbert, I.~Sergey. Compositional Verification of Composite Byzantine Protocols. CCS 2024. doi:10.1145/3658644.3690355.
\bibitem{flp} M.~J.~Fischer, N.~A.~Lynch, M.~S.~Paterson. Impossibility of Distributed Consensus with One Faulty Process. JACM 32(2), 1985.
\bibitem{pbft} M.~Castro, B.~Liskov. Practical Byzantine Fault Tolerance. OSDI 1999.
\end{thebibliography}

\end{document}
